\documentclass{situationist-preprint}

\title{Double Sensation and Its Asymmetry}
\msid{PHE-08-b}
\msfamily{PHE-08}
\msdiscipline{Phenomenology}
\mstype{Consequence}
\mskeywords{lived body; double sensation; reversibility; touch; zero-point}

\begin{document}
\maketitle

\begin{abstract}
PHE-08-a argues that the lived body is not an object, since the zero-point of orientation cannot be among the things it orients, and that in double sensation the touching and touched roles never coincide. We derive three consequences. First, the asymmetry of double sensation is not a limitation of attention but follows from the zero-point structure. Second, in each phase of double sensation exactly the touched hand is given as an object, and the touching hand is given otherwise. Third, double sensation differs in kind from touching another person's hand, so that it cannot serve as a model of intersubjective symmetry, nor that symmetry as a model of it.
\end{abstract}

\section{The asymmetry is structural}

When the right hand touches the left, Husserl observes, the left hand is felt as a thing with its smoothness and warmth, and also as the place of sensations \citep{Husserl1952}. Merleau-Ponty adds that the two hands never touch and are touched at once: one can pass from one role to the other, but the reversibility is always imminent and never realized \citep{MerleauPonty1964}. PHE-08-a states this as Proposition~4.2: at any moment the two roles of Definition~4.1 are not both played by each hand with respect to the other, and reversal is a succession, not a coincidence.

One might think the asymmetry merely reflects the narrowness of attention: one attends to the touching hand or to the touched one, and with better attention one could hold both. The first consequence denies this.

\begin{proposition}\label{p:struct}
The asymmetry of double sensation does not depend on the distribution of attention.
\end{proposition}

\begin{proof}
For $A$ to play the touching role is for $B$ to be given as an object through $A$'s sensations (Definition~4.1 of PHE-08-a), and so as located relative to a here (Definition~2.1). In that experience the here from which $B$ is located is the touching $A$, which is functioning as the zero-point. By Lemma~3.1 of PHE-08-a, the zero-point is not given, at the time at which it orients, as located in the space it orients. So $A$ is not given as object in that experience, whatever is attended to. The asymmetry follows from what the roles are, not from how attention is spread over them.
\end{proof}

\section{What is given as object}

\begin{corollary}\label{c:exact}
In each phase of double sensation, the touched hand is given as an object, and the touching hand is not; the touching hand is given as the place from which the touched hand is felt.
\end{corollary}

\begin{proof}
The first part is the definition of the touched role. The second follows from Proposition~\ref{p:struct}. That the touching hand is nonetheless given, and not absent, is the content of Husserl's localized sensations: the touching hand is felt as the bearer of the sensations through which the other is felt, which is to be given as functioning, not as located object.
\end{proof}

Corollary~\ref{c:exact} gives a precise sense to the claim of Theorem~3.2 of PHE-08-a that the body can be given as an object in part but not fully. In double sensation the part given as object and the part not so given are exactly the touched and the touching hand. When the roles reverse, what was touching is now given as object, but the earlier touching is retained as having been the touching, not as having been touched: retention does not convert the zero-point into an object after the fact. This is why the series of reversals does not add up to full givenness, as PHE-08-a argues in reply to its Objection~2.

\section{Not a model of the other}

\begin{corollary}\label{c:other}
Double sensation differs in kind from touching another's hand.
\end{corollary}

\begin{proof}
In a handshake each person's hand plays the touching role for that person, and both are touching at once, each from its own here. The symmetry is real because there are two zero-points. In double sensation there is one zero-point, so by Proposition~\ref{p:struct} at most one hand touches at a time.
\end{proof}

The corollary blocks two inferences. One cannot argue from the symmetry of the handshake that a perfectly attentive self-touching would be symmetrical, since the symmetry of the handshake depends on two orienting bodies. Nor can one argue, as some readings of Merleau-Ponty's late work suggest, that intersubjective contact is a double sensation writ large, so that it inherits its incompleteness; the incompleteness of double sensation comes from the single zero-point, which the handshake does not have.

\begin{objection}
Proprioception gives both hands at once, and gives them as located: I feel where both hands are while they touch.
\end{objection}

\begin{reply}
What is felt as located in such an experience is felt through the body image; the body schema, the system of sensorimotor capacities that orients the touching, functions without being given as located in that act \citep{Gallagher2005}. This is the distinction that PHE-08-a draws in reply to its Objection~3. The proprioceptive location of both hands therefore does not give the touching hand as object in its touching, and the asymmetry of the tactile roles is intact.
\end{reply}

\section{Conclusion}

The asymmetry of double sensation is structural. In each phase exactly one hand is given as object, and the other is given as the here from which it is felt; reversal never makes the two coincide; and the handshake, with its two zero-points, is a different structure. Double sensation shows the lived body given to itself as object in part, and never fully, as PHE-08-a concludes.

\begin{thebibliography}{99}\small
\bibitem[Gallagher(2005)]{Gallagher2005} Gallagher, S. (2005). \emph{How the Body Shapes the Mind}. Oxford University Press.
\bibitem[Husserl(1952)]{Husserl1952} Husserl, E. (1952). \emph{Ideen zu einer reinen Phänomenologie und phänomenologischen Philosophie. Zweites Buch}. Ed. M.~Biemel. Husserliana IV. Martinus Nijhoff.
\bibitem[Merleau-Ponty(1964)]{MerleauPonty1964} Merleau-Ponty, M. (1964). \emph{Le visible et l'invisible}. Ed. C.~Lefort. Gallimard.
\end{thebibliography}

\end{document}
